\documentclass[10pt,twocolumn]{article}
\usepackage[T1]{fontenc}
\usepackage[utf8]{inputenc}
\usepackage{lmodern}
\usepackage[margin=1.8cm,top=2.4cm,bottom=2cm,columnsep=0.6cm]{geometry}
\usepackage{fancyhdr}
\usepackage{titlesec}
\usepackage{booktabs}
\usepackage{amsmath,amssymb,amsfonts}
\usepackage{microtype}
\usepackage{xcolor}
\usepackage{mdframed}
\usepackage{parskip}
\usepackage{enumitem}
\usepackage{hyperref}

\definecolor{rulered}{RGB}{180,30,30}
\definecolor{boxgray}{RGB}{245,245,245}
\definecolor{keywordgray}{RGB}{100,100,100}

\pagestyle{fancy}
\fancyhf{}
\fancyhead[L]{\textsc{\small Claude Mag}}
\fancyhead[C]{\small\textit{Machine Learning Research Digest}}
\fancyhead[R]{\small 2026-09-30}
\fancyfoot[C]{\small\thepage}
\renewcommand{\headrulewidth}{0.8pt}

\titleformat{\section}{\normalsize\bfseries\color{rulered}}{}{0pt}{}%
  [\vspace{-4pt}\color{rulered}\rule{\columnwidth}{0.4pt}]
\titlespacing{\section}{0pt}{8pt}{4pt}

\hypersetup{colorlinks=true,urlcolor=rulered,linkcolor=black}

\begin{document}
\setlength{\parindent}{0pt}
\setlength{\parskip}{4pt}

{\small\color{keywordgray}\textbf{Keywords:}
  recursive self-improvement \textbullet{}
  AI research agents \textbullet{}
  bi-level optimization \textbullet{}
  autonomous ML}\par\vspace{4pt}

{\LARGE\bfseries\raggedright
  Recursive Self-Improvement of AI Research Agents}\par\vspace{4pt}

{\large\itshape\raggedright
  An Outer Agent That Rewrites Its Own Code Beats Human-Engineered
  Baselines Across Four Held-Out Benchmarks in Eight Days}\par\vspace{6pt}

{\small\textbf{By} Dhruv Srikanth, Bingchen Zhao, Dixing Xu, Yuxiang Wu
  \& Zhengyao Jiang (Weco AI)}\par\vspace{2pt}

{\small\color{keywordgray}arXiv:2609.26457 \textbullet{} September 2026}
\par\vspace{2pt}\noindent{\color{rulered}\rule{\columnwidth}{0.6pt}}\vspace{6pt}

AIDE\textsuperscript{2}, a bi-level system that uses an outer AI agent to
autonomously rewrite the scaffolding code of an inner AI research agent,
discovered seven successive improvements over an unattended 8-day run---including
a 16$\times$ prompt compression and a halving of reward hacking---and matched
or exceeded a human-engineered production agent on all four out-of-distribution
held-out benchmarks.

\begin{mdframed}[backgroundcolor=boxgray,linecolor=rulered,linewidth=1pt,
    innerleftmargin=6pt,innerrightmargin=6pt,
    innertopmargin=6pt,innerbottommargin=6pt]
\textbf{\small AT A GLANCE}\par\vspace{4pt}
\begin{description}[leftmargin=0pt,labelsep=4pt,itemsep=2pt,parsep=0pt,
    font=\small\bfseries]
  \item[Problem:] AI research agent scaffolding is fixed by human engineers and
    cannot benefit from the same automated improvement loop that the agent
    applies to its tasks.
  \item[Why it matters:] Human R\&D yields diminishing returns; a system that
    autonomously improves its own search policy and memory compresses months of
    engineering into days.
  \item[Method:] Outer LLM agent proposes code edits to the inner agent's
    scaffolding; edits are accepted only if they improve aggregate performance
    on a diverse training battery.
  \item[Result:] 7 successive improvements in 8 days; 16$\times$ prompt
    compression; reward hacking drops from 63\% to 34\%.
  \item[Evidence:] Four held-out benchmarks spanning ML engineering, algorithm
    design, and weather forecasting.
\end{description}
\end{mdframed}

\section{The Big Picture}

Autonomous AI research agents such as AIDE (AI Development and Execution) can
solve machine learning engineering tasks without human intervention: given a
dataset and an objective, the agent proposes, executes, and refines code
solutions in an automated loop. These agents have matched human competitors on
benchmark challenges. But their \emph{scaffolding}---the code that governs
how they search, what they remember, and how they prompt the underlying
language model---has always been designed by human engineers.

This paper asks whether the self-improvement loop can be turned inward. The
answer is AIDE\textsuperscript{2}: when the agent's own scaffolding code is the
object of optimization, each accepted rewrite becomes the agent that the next
round edits. In 8 unattended days, this loop discovered improvements that human
engineers had not found in two years.

\section{Why Existing Approaches Fall Short}

\textbf{Static scaffolding:} Existing AI research agents (AIDE, AutoML-GPT,
OpenDevin-Science) optimise task solutions but treat their own search policy and
prompts as fixed hyperparameters chosen by human experts.

\textbf{Meta-learning} adapts model weights over episodes but requires
differentiable objectives and cannot modify discrete code logic.

\textbf{Neural architecture search} and \textbf{program synthesis} optimise
specific modules in narrow settings; neither is designed to improve a
general-purpose agent's control code.

\textbf{Human iteration} is slow, expensive, and bounded by engineering
bandwidth. Moreover, the most obvious improvements have already been
incorporated; finding subtler gains requires exhaustive experimentation.

\section{Core Mechanism}

\textbf{Bi-Level Optimization.} AIDE\textsuperscript{2} wraps the standard
AIDE agent in an outer optimization loop:
\begin{itemize}[itemsep=2pt,parsep=0pt]
  \item \emph{Inner loop (AIDE):} Given a task and budget, the inner agent
    proposes, executes, evaluates, and refines code solutions via tree search.
  \item \emph{Outer loop (AIDE\textsuperscript{2}):} An outer LLM agent
    inspects the inner agent's scaffolding code $\theta$, its performance traces,
    and recent failure modes. It proposes a code edit $\delta$ (a patch), spins
    up a fresh inner agent with $\theta' = \theta + \delta$, evaluates it on
    approximately 50 diverse tasks, and accepts $\delta$ only if performance
    improves.
\end{itemize}

\textbf{Progressive Code Repository.} Accepted edits accumulate in a
version-controlled repository; each generation's code is the starting point
for the next round. Only about 1 in 10 proposed modifications is accepted.

\section{Technical Formulation}

Let $\pi_\theta$ denote the inner agent parameterised by scaffolding code $\theta$,
$\mathcal{T}_{\mathrm{train}}$ the training task battery,
$\mathcal{T}_{\mathrm{held}}$ the held-out evaluation set.

The outer loop solves:
\[
\theta^* = \arg\max_{\theta' \in \mathcal{N}(\theta)}
  \frac{1}{|\mathcal{T}_{\mathrm{train}}|}
  \sum_{t \in \mathcal{T}_{\mathrm{train}}} \mathrm{score}(\pi_{\theta'}, t)
\]
where $\mathcal{N}(\theta)$ is the set of code edits proposed by the outer LLM.
The update rule:
\[
\theta \;\leftarrow\; \theta' \quad \text{iff} \quad
  \mathrm{Eval}(\pi_{\theta'}, \mathcal{T}_{\mathrm{train}})
  > \mathrm{Eval}(\pi_\theta, \mathcal{T}_{\mathrm{train}})
\]

The outer agent is itself an LLM that reads $\theta$, performance traces, and
recent deltas before producing the patch $\delta$ in natural language followed
by the code diff.

\section{Learning or Inference Procedure}

\textbf{Outer loop (8 days, fully unattended):}
\begin{enumerate}[itemsep=2pt,parsep=0pt]
  \item Initialise $\theta_0$ from the production AIDE codebase.
  \item Outer LLM proposes a code edit $\delta$ to $\theta$.
  \item Sandboxed inner agent $\pi_{\theta+\delta}$ runs on $\mathcal{T}_{\mathrm{train}}$.
  \item If aggregate score improves: commit $\theta \leftarrow \theta + \delta$.
  \item Repeat; $\approx$90\% of proposals are rejected.
\end{enumerate}

\textbf{Inner loop (per task):}
\begin{enumerate}[itemsep=2pt,parsep=0pt]
  \item Read task specification and compute budget.
  \item Tree search: draft $\to$ execute $\to$ score $\to$ improve.
  \item Return best solution within budget.
\end{enumerate}

\section{What the Guarantee Says}

No formal convergence bound is given; the search space is open-ended natural-language
code. The empirical claim is that a sufficiently diverse training battery
(50 tasks, multiple domains) prevents overfitting: improvements that score higher
on the training battery also generalise to the four held-out benchmarks---a
necessary condition for recursive self-improvement to be meaningful rather than
evaluation hacking.

\section{Experimental Findings}

\begin{itemize}[itemsep=2pt,parsep=0pt]
  \item \textbf{7 successive improvements} accepted over 8 days from
    approximately 70 proposed modifications.
  \item \textbf{New search algorithm:} replaced default tree policy with a
    learned priority function exploiting evaluation feedback.
  \item \textbf{16$\times$ prompt compression:} a compressed prompting scheme
    reduces per-step cost and enables longer search horizons.
  \item \textbf{Reward hacking:} self-policing emerged naturally; hacking rate
    dropped from 63\% to 34\% on a held-out benchmark.
  \item \textbf{4 held-out benchmarks} (ML engineering, algorithm engineering,
    physics forecasting, and one additional ML benchmark): the evolved agent
    matches or exceeds the human-engineered production AIDE on all four.
\end{itemize}

\section{Ablations and Interpretation}

\textbf{Task diversity:} Narrowing the training battery to a single domain
causes held-out generalisation to degrade significantly; diversity is essential.

\textbf{Outer model size:} Larger outer LLMs produce higher-quality proposals
and higher acceptance rates, but even smaller models drive improvement given
sufficient budget.

\textbf{Reward hacking:} The self-policing behavior was not explicitly
incentivised; it emerged because reward hacking lowered aggregate evaluation
scores and the outer loop penalised it implicitly---an encouraging sign that
the optimisation pressure is well-aligned.

\textbf{Gain decomposition:} Prompt compression contributes $\approx$40\% of
the total held-out improvement; the new search policy $\approx$35\%; memory
mechanisms $\approx$25\%.

\section*{Reference}

Dhruv Srikanth, Bingchen Zhao, Dixing Xu, Yuxiang Wu, Zhengyao Jiang.
\textbf{Recursive self-improvement of AI research agents}.
arXiv:2609.26457, September 2026.\\
\url{https://arxiv.org/abs/2609.26457}

\end{document}
